% !TEX program = xelatex
% PyQCD 与 donghx Gradient_Flow_GluonPDF 算法审计、修正与对照报告
\documentclass[UTF8,a4paper,11pt]{ctexart}

\usepackage[a4paper,margin=2.05cm,headheight=14pt]{geometry}
\usepackage{amsmath,amssymb,mathtools,bm}
\usepackage{booktabs,longtable,array,tabularx,multirow}
\usepackage{graphicx}
\usepackage{xcolor}
\usepackage{listings}
\usepackage{tikz}
\usetikzlibrary{arrows.meta,positioning,calc,fit,shapes.geometric}
\usepackage{url}
\usepackage{hyperref}

\definecolor{pyblue}{RGB}{36,91,163}
\definecolor{qudared}{RGB}{190,65,55}
\definecolor{stagegray}{RGB}{90,105,115}
\definecolor{goodgreen}{RGB}{30,132,73}
\definecolor{warnorange}{RGB}{197,90,17}
\hypersetup{
  colorlinks=true,linkcolor=pyblue,urlcolor=pyblue,citecolor=pyblue,
  pdftitle={PyQCD 与 donghx 梯度流胶子 OPE 算法审计与修正}
}
\urlstyle{tt}
\newcommand{\pathref}[1]{\path{#1}}
\newcommand{\code}[1]{\texttt{\detokenize{#1}}}
\newcommand{\codeurl}[1]{\nolinkurl{#1}}
\newcommand{\dd}{\mathrm{d}}
\newcommand{\tr}{\operatorname{tr}}
\newcommand{\Tr}{\operatorname{Tr}}
\newcommand{\C}{\mathbb C}
\newcommand{\R}{\mathbb R}
\newcommand{\SU}{\mathrm{SU}}
\newcommand{\MSbar}{\overline{\mathrm{MS}}}
\newcommand{\GF}{\mathrm{GF}}
\newcommand{\odd}{\mathrm{o}}
\newcommand{\even}{\mathrm{e}}
\newcommand{\verified}{\textcolor{goodgreen}{\textbf{确证}}}
\newcommand{\inferred}{\textcolor{pyblue}{\textbf{推导}}}
\newcommand{\unverified}{\textcolor{warnorange}{\textbf{未验证}}}
\setlength{\parindent}{2em}
\setlength{\parskip}{0.35em}
\emergencystretch=3em
\setlength{\tabcolsep}{4pt}
\renewcommand{\arraystretch}{1.08}

\lstset{
  basicstyle=\ttfamily\small,
  backgroundcolor=\color{black!3},
  frame=single,
  breaklines=true,
  breakatwhitespace=false,
  keepspaces=true,
  columns=flexible,
  keywordstyle=\color{pyblue}\bfseries,
  commentstyle=\color{goodgreen!70!black},
  stringstyle=\color{pyblue}
}

\title{PyQCD 与 donghx 梯度流胶子 OPE/GluonPDF 链\\
算法审计、物理解析、实现修正与对照测试}
\author{PyQCD 代码—物理分析记录}
\date{2026 年 9 月 9 日}

\begin{document}
\maketitle

\begin{abstract}
本文审计并修正 \pathref{/root/PyQCD} 与
\pathref{/root/PyQCD/refer/donghx/Gradient_Flow_GluonPDF} 中的梯度流胶子
算法。工作重点是薄链接输入、四维 Wilson flow、四叶 Clover 场强、直线
Wilson 线胶子 OPE、schema-v2 原始正负 $z$ 分支、核子两点函数驱动的
disconnected $C_3/C_2$ 统计，以及一圈梯度流到
$\MSbar$ 准算符的端点转换。报告首先给出参考实现的物理正确性与证据等级，
再说明既有 PyQCD legacy/TMD 路径为何不能直接冒充该参考链，最后给出
新增模块、严格验证器和独立契约测试。

本轮的核心结果是：OPE 组件、Wilson 线几何、复数 helicity
$\texttt{pol35}$ 协方差估计、共同 $\texttt{nopol}$ 分母以及
delete-one jackknife 已在算法级复现；参考 ratio 估计器的
$C_3$、$C_2$、ratio、jackknife 均值、偏差修正误差和有效掩码逐项达到
$\max|{\rm diff}|=0$。仓库全量测试为 $106$ 项通过。由于当前工作树没有
参考目录中的完整 $L32\times96$、$231$ 组态 OPE/2pt 生产产物，本文不把
算法级复现夸大为真实 $231$ 组态逐元素生产复现；有限流 bare observable
也不被称为已经完成的物理 gluon PDF。
\end{abstract}

\tableofcontents

\section{范围、结论与证据等级}

\subsection{任务对象与判断标准}

本报告的工作对象是 PyQCD 中与 Gradient Flow/GluonPDF 直接相关的部分，
对照对象是
\pathref{/root/PyQCD/refer/donghx/Gradient_Flow_GluonPDF} 的全部算法说明、
生产脚本、验证脚本和测试脚本。参考目录的核心文件包括：

\begin{itemize}
  \item \pathref{gradient-flow-gluon-ratio.md}：流、Clover、Wilson 线、
        schema-v2 与 ratio 约定；
  \item \pathref{calc_ratio_gradient_flow.py}：完整复数 disconnected
        估计器及 delete-one jackknife；
  \item \pathref{validate_output_v5.py} 与 \pathref{validate_ratio.py}：
        产物形状、线性恒等式、有限性和掩码守卫；
  \item \pathref{gluon_flow_to_quasi.py}：一圈 $\GF\to\MSbar$ 准算符端点
        转换；
  \item \pathref{helicity_operator_comparison.py} 与
        \pathref{test_estimators.py}：helicity 物理选择和统计 oracle。
\end{itemize}

参考目录声称生产实现位于 \code{src/flowed_gluon_ope.cpp}，但该源码和完整
生产二进制不在当前对照目录中。因此，关于 OPE 数值语义的直接依据是
Markdown、Python 生产/验证代码和小数组代数测试；不存在的 C++ 源码不被
虚构为已审阅证据。

\begin{table}[htbp]
\centering
\caption{本报告的证据等级}
\begin{tabularx}{\textwidth}{@{}p{0.20\textwidth}p{0.25\textwidth}X@{}}
\toprule
等级 & 含义 & 本轮实例 \\
\midrule
源码确证 &
可由参考脚本、PyQCD 函数和测试代码直接定位 &
流步长、Clover 归一、OPE 轴、ratio 公式、jackknife 展开 \\
\midrule
代数/物理推导 &
由链路定义、规范变换和线性关系推出 &
Wilson 线规范不变性、$z=0$ odd 为零、组合分量重构 \\
\midrule
独立数值 oracle &
同一输入数组经参考算法与 PyQCD 算法分别计算 &
ratio 六类数值以及有效掩码逐项 $\max|{\rm diff}|=0$ \\
\midrule
真实生产对照 &
同一完整组态、同一 manifest、同一产物的逐元素比较 &
当前缺少参考 $L32\times96$/231 组态产物，不能宣称已完成 \\
\bottomrule
\end{tabularx}
\end{table}

\subsection{结论先行}

\begin{enumerate}
  \item 参考链的受控终点是\textbf{有限流时间 bare disconnected
        ratio}，而不是已经重整化、匹配并完成连续极限的 gluon PDF。
  \item PyQCD 新增 schema-v2 层，保留旧的 legacy/TMD 默认语义；这样可以
        同时复现参考约定和保护已经验证的 docker 管线。
  \item OPE 的生产投影固定为 color-traceless Clover；legacy
        untraced 投影仅作为有限格距诊断和兼容层保存。
  \item helicity ratio 必须在协方差构造时保留完整复数
        $\texttt{pol35}$，而且 numerator 使用 $\texttt{pol35}$、
        denominator 使用共同的 $\texttt{nopol}$。最后才作
        $\operatorname{Re}[-iR_H]=\operatorname{Im}R_H$ 投影。
  \item 当前已完成的是算法、接口、验证和小数组 oracle 闭环；真实
        $N_{\rm conf}=231$ 产物、统计外推和最终物理 PDF 仍需补充数据与
        多尺度控制。
\end{enumerate}

\section{统一约定与数据边界}

\subsection{格点、方向和量纲}

规范链接数组采用
\begin{equation}
 U\;:\;(N_t,N_z,N_y,N_x,4,N_c,N_c),
\qquad N_c=3,
\label{eq:gauge-layout}
\end{equation}
其数组轴顺序是 $(t,z,y,x)$，而 Lorentz/link 方向标签是
\[
\mu=0:x,\qquad \mu=1:y,\qquad \mu=2:z,\qquad \mu=3:t.
\]
因此一个 Lorentz 方向 $\mu$ 对应数组空间轴 $3-\mu$。所有空间和使用
周期边界；代码中的 \code{roll} 正是
$f(x+\hat\mu)$ 或 $f(x-\hat\mu)$ 的离散实现。

参考生产固定
\begin{equation}
 z_{\rm dir}=2,\qquad
 \tau=\frac{t}{a^2},\qquad
 \epsilon=0.01,\qquad
 r_{\rm flow}/a=\sqrt{8\tau}.
\label{eq:flow-units}
\end{equation}
$\tau$ 是无量纲流时间；只有在引入格距 $a$ 后才可换成物理流时间
$t_{\rm phys}=\tau a^2$。如果 $a$ 以 fm 给出，则
\begin{equation}
 t_{\rm phys}[{\rm GeV}^{-2}]
 =\tau\left(a[{\rm fm}]\times 5.067730716\right)^2.
\label{eq:flow-convert}
\end{equation}
混用 $\tau$、$t$ 和 $t_{\rm phys}$ 会直接改变小流时间匹配系数，故新模块将
它们分别写入 metadata。

\subsection{参考 OPE schema-v2}

单个组态的 OPE 数据轴为
\begin{equation}
(\texttt{field\_projection},\texttt{channel},\texttt{z\_orientation},
\texttt{component},z,t)
=(2,2,4,6,N_z,N_t).
\label{eq:ope-shape}
\end{equation}
固定标签如下：

\begin{table}[htbp]
\centering
\caption{schema-v2 的轴标签}
\begin{tabularx}{\textwidth}{@{}p{0.25\textwidth}p{0.30\textwidth}X@{}}
\toprule
轴 & 标签 & 物理/数值意义 \\
\midrule
field projection &
\texttt{legacy\_untraced}, \texttt{traceless} &
有限格距未去迹诊断与生产用 $\SU(3)$ 无迹投影 \\
\midrule
channel &
\texttt{unpolarized}, \texttt{helicity} &
第二个场强分别使用 $F$ 或 $\widetilde F$ \\
\midrule
z orientation &
\texttt{plus\_z\_raw}, \texttt{minus\_z\_raw},
\texttt{even\_sum}, \texttt{odd\_difference} &
保存原始方向，even/odd 不除以 $2$ \\
\midrule
component &
\texttt{combined}, \texttt{Mtiti}, \texttt{Mijij},
\texttt{ti}, \texttt{tj}, \texttt{ij\_single} &
保存独立分量和组合分量，便于恒等式与 helicity 诊断 \\
\bottomrule
\end{tabularx}
\end{table}

新产物状态明确记录为
\code{flowed_bare_ope_observable}。它不隐含 $Z_R$、混合、Collins--Soper
演化、软因子、LaMET matching 或连续极限。

\subsection{工作树中的真实数据覆盖}

当前默认数据目录 \pathref{/root/PyQCD/data} 的实际覆盖为：

\begin{table}[htbp]
\centering
\caption{当前可访问数据与缺口}
\begin{tabularx}{\textwidth}{@{}p{0.32\textwidth}p{0.23\textwidth}X@{}}
\toprule
路径 & 已有内容 & 对本任务的意义 \\
\midrule
\pathref{data/cmp1\_4150/} &
组态 4150 的 VVV 数组 &
可用于既有蒸馏低层对照，不能替代流化 OPE/2pt 生产数据 \\
\midrule
\pathref{data/cmp1\_cache/gauge\_4150.npy} &
少量规范场缓存 &
可做小规模 OPE/规范协变测试；不构成共同 231 组态 manifest \\
\midrule
参考 $L32\times96$ 生产根 &
当前工作树中不存在 &
不能做真实 OPE、C2、ratio 的逐文件或逐元素对照 \\
\bottomrule
\end{tabularx}
\end{table}

\section{参考算法的物理解析}

\subsection{Wilson flow：耗散的规范场演化}

连续梯度流方程为
\begin{equation}
 \partial_t B_\mu(t,x)=D_\nu G_{\nu\mu}(t,x),
 \qquad B_\mu(0,x)=A_\mu(x).
\label{eq:continuum-flow}
\end{equation}
它把四维规范场沿作用量下降方向平滑；在正的流时间，半径约为
$\sqrt{8t}$ 的紫外涨落被压制。格点上用链接 $V_\mu(t,x)$ 表示流场，
Wilson 作用量的六个 staple 给出
\begin{equation}
 \Omega_\mu(x)=\sum_{\nu\ne\mu}
 \left[
 U_\nu(x)U_\mu(x+\hat\nu)U_\nu^\dagger(x+\hat\mu)
 U_\nu^\dagger(x-\hat\nu)U_\mu(x-\hat\nu)
 U_\nu(x-\hat\nu+\hat\mu)
 \right].
\label{eq:staple}
\end{equation}
实际代码等价地把每一项按矩阵乘法写成正、负两个三段路径。

流右端首先形成 $\Omega_\mu V_\mu^\dagger$，再取无迹反厄米部分：
\begin{equation}
 Z_\mu(V)=P_{\rm ah}\!\left[\Omega_\mu V_\mu^\dagger\right],
\qquad
 P_{\rm ah}(X)
 =\frac{X-X^\dagger}{2}
 -\frac{\tr(X-X^\dagger)}{2N_c}\mathbf 1.
\label{eq:flow-generator}
\end{equation}
因此 $Z_\mu\in\mathfrak{su}(3)$。若指数精确计算，则
$\exp(\epsilon Z_\mu)\in\SU(3)$；实现中的 \code{su3\_exp} 采用小步长四阶
截断并做 determinant 归一，随后以幺正性和 determinant 误差作守卫。

\subsubsection{三阶 Runge--Kutta}

Lüscher 三阶更新写成
\begin{align}
 W_0&=V_t,\\
 W_1&=\exp\left(\frac{\epsilon}{4}Z_0\right)W_0,\\
 W_2&=\exp\left(\frac{8\epsilon}{9}Z_1
              -\frac{17\epsilon}{36}Z_0\right)W_1,\\
 V_{t+\epsilon}
 &=\exp\left(\frac{3\epsilon}{4}Z_2
              -\frac{8\epsilon}{9}Z_1
              +\frac{17\epsilon}{36}Z_0\right)W_2.
\label{eq:rk3}
\end{align}
当只给出 $\tau$ 和 $\epsilon$ 时，步数为
\begin{equation}
 N_{\rm step}=\left\lceil\frac{\tau}{\epsilon}\right\rceil,
\qquad \delta t=\frac{\tau}{N_{\rm step}}.
\label{eq:flow-steps}
\end{equation}
最后一步使用 $\delta t$ 而非盲目使用 $\epsilon$，故实际流时间恰好命中目标
$\tau$。显式传入 $N_{\rm step}$ 时，metadata 同时保存用户步数和真实步长。

物理上应监测 Wilson plaquette 作用量密度
\begin{equation}
 s_W(x)=\frac{1}{6}\sum_{\mu<\nu}
 \left(1-\frac{\operatorname{ReTr}P_{\mu\nu}(x)}{N_c}\right),
\label{eq:wilson-density}
\end{equation}
而不应把任意 Clover $G^2$ 离散量误当作严格逐点单调量。新 OPE metadata
保存流前、流后的平均 plaquette、最大链接幺正性误差和最大
$|\det U-1|$，验证器要求有限、流后 plaquette 不下降并满足有限阈值。

\subsection{Clover 场强、无迹投影与对偶}

四叶 Clover 的基本矩阵 $Q_{\mu\nu}(x)$ 是同一中心四个有向 plaquette
的和，场强按参考约定为
\begin{equation}
 F_{\mu\nu}(x)
 =-\frac{i}{8}\left[Q_{\mu\nu}(x)-Q_{\mu\nu}^\dagger(x)\right].
\label{eq:clover}
\end{equation}
该对象是有限格距的 color 矩阵；由于四叶和与高阶格距项，未投影时可以含有
单位阵分量。生产投影为
\begin{equation}
 F_{\mu\nu}^{\rm tl}
 =F_{\mu\nu}
 -\frac{\tr F_{\mu\nu}}{N_c}\mathbf 1,
\qquad
 F_{\nu\mu}=-F_{\mu\nu}.
\label{eq:traceless}
\end{equation}
PyQCD 同时保存 \code{legacy\_untraced} 和 \code{traceless}，但只把后者
作为 production selector；两者的差异是有限格距离散差异，不能在分析中静默
混用。

对偶场强由 Euclidean Levi--Civita 张量构造：
\begin{equation}
 \widetilde F_{\mu\nu}
 =\frac{1}{2}\epsilon_{\mu\nu\rho\sigma}F_{\rho\sigma}.
\label{eq:dual}
\end{equation}
代码只存六个 canonical Lorentz pair，反向请求通过反对称性取负；这避免了
重复计算 Clover，并使 $F_{\mu\nu}$ 和 $\widetilde F_{\mu\nu}$ 的轴顺序明确。

\subsection{直线 Wilson 线与规范不变性}

沿 $z_{\rm dir}$ 的正方向双局域量为
\begin{equation}
 M_{\mu\lambda;\nu\rho}^{+}(z)
 =\sum_{\boldsymbol x}
 \Tr\left[
 F_{\mu\lambda}(x+z\hat z)\,
 W^\dagger(x+z,x)\,
 F_{\nu\rho}(x)\,
 W(x,x+z)
 \right].
\label{eq:bilocal-plus}
\end{equation}
负方向使用 $x-z\hat z$ 和从 $x-z$ 到 $x$ 的正向路径：
\begin{equation}
 M_{\mu\lambda;\nu\rho}^{-}(z)
 =\sum_{\boldsymbol x}
 \Tr\left[
 F_{\mu\lambda}(x-z\hat z)\,
 W(x-z,x)\,
 F_{\nu\rho}(x)\,
 W^\dagger(x,x-z)
 \right].
\label{eq:bilocal-minus}
\end{equation}
在 $z=0$ 两式退化为同一个局域矩阵乘积，这是
\begin{equation}
 M^+(0)=M^-(0),\qquad
 M_{\rm odd}(0)=M^+(0)-M^-(0)=0
\label{eq:z0-limit}
\end{equation}
的几何原因。

规范变换下
\begin{equation}
 F(x)\mapsto G(x)F(x)G^\dagger(x),\qquad
 W(x,x+z)\mapsto G(x)W(x,x+z)G^\dagger(x+z).
\label{eq:gauge-transform}
\end{equation}
代入式 \eqref{eq:bilocal-plus} 后，端点的 $G$ 和 $G^\dagger$ 在 trace
中相消，所以完整空间和是规范不变的。只在某个点比较矩阵元素、或把两端
场强放在不同的链接几何上而没有 Wilson 线，都会破坏这个论证。

\subsection{六个 OPE 组件与物理选择}

令 Wilson 线沿 $z$，横向方向为 $i,j$，Euclidean 时间方向为 $t=3$。
定义三个 primitive：
\begin{align}
 T_i&=M_{ti;ti},&
 T_j&=M_{tj;tj},&
 S_0&=M_{ij;ij}.
\end{align}
保存组件为
\begin{align}
 \texttt{ti}&=T_i,&
 \texttt{tj}&=T_j,&
 \texttt{ij\_single}&=S_0,\\
 M_{titi}&=T_i+T_j,&
 M_{ijij}&=2S_0.
\label{eq:components}
\end{align}
因此 \code{Mtiti} 和 \code{Mijij} 是可由 primitive 逐点重构的独立诊断。

无极化通道使用
\begin{equation}
 O_U=M_{titi}-M_{ijij}
 =M_{tx;tx}+M_{ty;ty}-2M_{xy;xy}.
\label{eq:unpol-combination}
\end{equation}
helicity 通道将第二端点替换为对偶场强，参考生产数组的
\code{combined} 保存的是
\begin{equation}
 O_H^{\rm stored}=M_{titi}+M_{ijij}.
\label{eq:hel-stored}
\end{equation}
在 helicity 算符比较中，真正的理论目标由原始方向的
\begin{equation}
 O_H^{\rm physical}
 =\left(M_{titi}-M_{ijij}\right)_{\odd}
\label{eq:hel-physical}
\end{equation}
显式 selector 得到；$O_H^{\rm stored}$ 保留作诊断。把保存约定的加号直接
称为唯一物理 helicity 定义，会把诊断和目标混淆。

方向投影不引入 $1/2$：
\begin{equation}
 O_{\even}=O_{+z}+O_{-z},\qquad
 O_{\odd}=O_{+z}-O_{-z}.
\label{eq:even-odd}
\end{equation}
因此验证器可以同时检查方向线性关系、$z=0$ odd 零值和六组件组合恒等式。

\section{disconnected $C_3/C_2$ 统计}

\subsection{输入轴与 numerator/denominator}

新 ratio 入口
\code{calculate_gradient_flow_ratios}（同时提供参考同名别名
\code{calculate_ratios}）接收
\begin{align}
 O&:(N_{\rm conf},2,4,N_{\rm comp},N_z,N_t),\\
 C_2&:(2,N_{\rm dir},N_{\rm conf},N_{\rm tsep},N_t,N_p).
\label{eq:ratio-input}
\end{align}
channel $0$ 是 \code{nopol}，channel $1$ 是 \code{pol35}。对任意 channel
$\chi$，C3 的 numerator 使用对应的投影：
\begin{equation}
 C_{3,\chi}
 =\left\langle O_\chi C_{2,\chi}\right\rangle_{\rm conf,src}
 -\left\langle O_\chi\right\rangle_{\rm conf,src}
  \left\langle C_{2,\chi}\right\rangle_{\rm conf,src}.
\label{eq:c3-cov}
\end{equation}
但 ratio 的共同分母固定为
\begin{equation}
 C_{2,\rm den}
 =\left\langle C_{2,\rm nopol}\right\rangle_{\rm conf,src},
\qquad
 R_\chi=\frac{C_{3,\chi}}{C_{2,\rm den}}.
\label{eq:common-denominator}
\end{equation}
这一区分是本轮修正的核心之一：helicity 的 pol35 只进入 numerator，
不能再用 ensemble mean pol35 作为近零分母。

插入时间 $\tau_{\rm ins}$ 通过周期源移位实现：
\begin{equation}
 O(t_{\rm src}+\tau_{\rm ins})
 =\operatorname{take}\left[
 O,\,(t_{\rm src}+\tau_{\rm ins})\bmod N_t
\right].
\label{eq:source-shift}
\end{equation}
只有 $0\leq\tau_{\rm ins}\leq t_{\rm sep}$ 的位置有效，输出共同轴上的
其余位置全部填为 NaN，而不是静默填零。

\subsection{复数 helicity 相位}

若
\begin{equation}
 R_H=a+ib,
\end{equation}
Euclidean helicity 的最终投影为
\begin{equation}
 \operatorname{Re}[-iR_H]=b=\operatorname{Im}R_H.
\label{eq:hel-phase}
\end{equation}
但是这个投影只能发生在配置协方差、真空扣除、ratio 和 jackknife 完成之后。
若先把 $\texttt{pol35}$ 截成实部或虚部，则
\[
\langle O\,\operatorname{Im}C_2\rangle
-\langle O\rangle\langle\operatorname{Im}C_2\rangle
\]
已经不是完整复数协方差，统计量改变且无法由最终取虚部补回。

\subsection{delete-one jackknife 的代数展开}

对一个固定 channel、方向、$z$、组件、separation、插入时间和动量，设
$N=N_{\rm conf}$。定义每个组态的 source 平均量
\begin{equation}
 a_k=\langle O_kC_k\rangle_s,\qquad
 o_k=\langle O_k\rangle_s,\qquad
 c_k=\langle C_k\rangle_s,
\end{equation}
以及全样本平均
\begin{equation}
 A=\frac1N\sum_k a_k,\qquad
 O=\frac1N\sum_k o_k,\qquad
 C=\frac1N\sum_k c_k.
\end{equation}
删除第 $k$ 个组态后，
\begin{equation}
 C_{3,(-k)}
 =\frac{NA-a_k}{N-1}
 -\frac{(NO-o_k)(NC-c_k)}{(N-1)^2}.
\label{eq:jk-c3}
\end{equation}
共同 nopol 分母同理：
\begin{equation}
 C_{2,(-k),\rm den}
 =\frac{ND-d_k}{N-1},
\qquad D=\frac1N\sum_k d_k.
\label{eq:jk-c2}
\end{equation}
然后在每一个删除样本内重新做
\begin{equation}
 R_{(-k)}=\frac{C_{3,(-k)}}{C_{2,(-k),\rm den}}.
\label{eq:jk-ratio}
\end{equation}
PyQCD 通过总和与 Einstein contraction 展开式计算式
\eqref{eq:jk-c3}，避免显式形成最大的
$(N,N_{\rm obs},N_t,N_p)$ 临时数组；代数上与逐个删除再直接计算完全等价。

jackknife 均值、偏差修正和实/虚误差分别为
\begin{align}
 \overline R_{\rm JK}
 &=\frac1N\sum_k R_{(-k)},\\
 R_{\rm bc}
 &=N R_{\rm central}-(N-1)\overline R_{\rm JK},\\
 \sigma_{\rm Re}
 &=\left[\frac{N-1}{N}\sum_k
   \left(\operatorname{Re}R_{(-k)}
        -\operatorname{Re}\overline R_{\rm JK}\right)^2\right]^{1/2},\\
 \sigma_{\rm Im}
 &=\left[\frac{N-1}{N}\sum_k
   \left(\operatorname{Im}R_{(-k)}
        -\operatorname{Im}\overline R_{\rm JK}\right)^2\right]^{1/2}.
\label{eq:jk-errors}
\end{align}
误差不把复数幅值与实虚部混成一个未定义的标量。

\section{一圈 $\GF\to\MSbar$ 准算符转换}

\subsection{端点系数与线性质量项}

参考 \pathref{gluon\_flow\_to\_quasi.py} 先把 $\tau$ 转为
$t[{\rm GeV}^{-2}]$，然后使用
\begin{align}
 c_\parallel&=1,\\
 c_\perp&=1+\frac{\alpha_s C_A}{4\pi}
 \log\left(2\mu^2t\,e^{\gamma_E}\right),\\
 \delta_m&=-\frac{\alpha_s C_A}{4\pi}
 \sqrt{\frac{2\pi}{t}}.
\label{eq:flow-quasi-coeff}
\end{align}
沿 Wilson 线长度 $z$ 的指数因子为
\begin{equation}
 L(z)=\exp[-\delta_m |z|],
\qquad z[{\rm GeV}^{-1}]
 =z/a\times a[{\rm fm}]\times5.067730716.
\label{eq:line-factor}
\end{equation}
于是两个通道的转换因子为
\begin{equation}
 Z_U(z)=\frac{L(z)}{c_\perp^2},
\qquad
 Z_H(z)=\frac{L(z)}{c_\parallel c_\perp}.
\label{eq:flow-quasi-factor}
\end{equation}
PyQCD 新增 \pathref{pyqcd/renorm/\_gradient\_flow\_quasi.py}，保留参考
\code{coefficients}、\code{match\_one} 的同名兼容入口，并提供语义化别名：
\begin{itemize}
  \item \codeurl{gradient_flow_to_quasi_coefficients}；
  \item \codeurl{match_gradient_flow_ope_to_quasi}。
\end{itemize}

\subsection{为什么必须重新构造 combined}

端点因子已经作用于 \code{Mtiti} 和 \code{Mijij}。若输入数组中的
\code{combined} 先被乘因子、再与匹配后的独立分量不一致地混用，则会留下
不受控的组合误差。新实现先按 channel 乘 $Z_U/Z_H$ 和 $L(z)$，再按
\begin{equation}
 \texttt{combined}_U
 =\texttt{Mtiti}_U-\texttt{Mijij}_U,\qquad
 \texttt{combined}_H
 =\texttt{Mtiti}_H+\texttt{Mijij}_H
\label{eq:matched-reconstruct}
\end{equation}
重构。该层仍只是准算符转换；光锥 matching、软因子、快度演化和连续极限
没有被这两个公式自动完成。

\section{既有 PyQCD 的偏差与修正策略}

\subsection{为什么不直接覆盖 legacy API}

PyQCD 原有 \pathref{pyqcd/operator/\_gluon\_ope.py} 是 docker 成功实例的
legacy 路径。它的 \code{plaquette\_clover} 和
\code{gluon\_ope\_operator\_z0} 对历史 finite-$a$、未去迹 OPE 有明确约定；
该路径已经被既有测试和成功实例使用。参考 schema-v2 则要求：

\begin{itemize}
  \item thin link 直接进入四维 Wilson flow；
  \item 同时输出未去迹和 traceless 两种 field projection；
  \item 保存 $\pm z$ 原始分支、even/odd 分支和六种组件；
  \item production analysis 固定选择 traceless；
  \item ratio 中 helicity 使用完整复数 pol35 numerator 和 nopol denominator。
\end{itemize}

如果直接改变 legacy 默认值，旧的 docker 一致性和既有 HYP/TMD 语义都会被
无关地改变。因此本轮采用独立模块
\pathref{pyqcd/operator/\_gradient\_flow\_gluon\_ope.py}，在 metadata 中
明确 schema 和 status；旧入口保持原语义，两个层不互称。

\subsection{问题—修正对照表}

\begin{longtable}{@{}p{0.22\textwidth}p{0.31\textwidth}p{0.37\textwidth}@{}}
\caption{工作对象与参考算法的关键偏差、修正和边界}\label{tab:fixes}\\
\toprule
既有风险/偏差 & 物理后果 & 本轮修正 \\
\midrule
\endfirsthead
\toprule
既有风险/偏差 & 物理后果 & 本轮修正 \\
\midrule
\endhead
把 legacy finite-$a$ 未去迹 OPE 当成 production traceless OPE &
有限格距单位阵分量污染，且无法与 reference schema 的 projection 轴对应 &
新增双 projection schema；生产选择 \code{traceless}，legacy 仅诊断 \\
\midrule
只保存一个正向 $z$ 或只保存已组合结果 &
无法检查 $z=0$ 极限、even/odd 方向关系和 helicity 选择敏感性 &
保存 \code{plus\_z\_raw}/\code{minus\_z\_raw}、四方向标签和六组件 \\
\midrule
helicity 先取 pol35 实部/虚部再做 covariance &
丢失复数配置相关性，所得量不是完整 disconnected 协方差 &
ratio estimator 全程保留 \code{complex128}，末端才取
$\operatorname{Re}[-iR_H]$ \\
\midrule
helicity ratio 用 pol35 ensemble mean 作分母 &
分母接近零，ratio 数值病态且与参考 v2 约定不符 &
channel numerator 分开取 nopol/pol35，两个 channel 共享 nopol 分母 \\
\midrule
删除组态后只重算 numerator 或只重算 C2 &
vacuum subtraction 与分母的 jackknife 相关性被破坏 &
按式 \eqref{eq:jk-c3}--\eqref{eq:jk-ratio} 在每个 resample 内同时重算 \\
\midrule
无效 insertion 用零填充 &
padding 被误认为物理数据，均值/误差被系统性拉偏 &
所有无效位置使用 NaN，并保存 \code{valid\_insertion\_mask} \\
\midrule
把现有 TMD staple/hybrid 路径当作 reference 直线 OPE &
HYP、$\epsilon$、staple 长度、横向分离和重整化方案被混合 &
新 schema-v2 只实现薄链接直线链；现有 TMD 默认路径保持隔离 \\
\midrule
一圈 GF-to-quasi 因子和现有 SFTX 系数同名混用 &
流时间方案、端点因子和算符层次混淆 &
新增独立 \code{\_gradient\_flow\_quasi.py}，只实现参考一圈端点转换 \\
\bottomrule
\end{longtable}

\subsection{严格校验器的必要性}

OPE 验证器
\code{validate\_gradient\_flow\_gluon\_ope} 检查：

\begin{itemize}
  \item schema、status、dtype、轴顺序、标签和 shape；
  \item 流前后 plaquette 有限、链接接近 $\SU(3)$ 且流方向正确；
  \item $z=0$ 正负方向相等、even/odd 线性恒等式；
  \item \code{Mtiti=ti+tj}、\code{Mijij=2*ij\_single} 以及两 channel
        combined 重构。
\end{itemize}

ratio 验证器
\code{validate\_gradient\_flow\_ratio\_results} 进一步检查：

\begin{itemize}
  \item 所有 central/jackknife/error 数组的轴和复数/实数 dtype；
  \item \code{c2\_mean} 形状、有限性、非零性和两个 channel 的严格共享；
  \item 有效 insertion 必须有限，无效 insertion 必须全为 NaN；
  \item central ratio 必须等于 \code{c3\_mean}/\code{c2\_mean[nopol]}；
  \item ratio 与 C3 的方向线性恒等式，以及 helicity odd 的 $z=0$ 零值。
\end{itemize}

\section{实现结构与调用链}

\subsection{模块映射}

\begin{table}[htbp]
\centering
\caption{最终 PyQCD 实现映射}
\scriptsize
\begin{tabularx}{\textwidth}{@{}p{0.35\textwidth}p{0.26\textwidth}X@{}}
\toprule
文件 & 入口 & 职责 \\
\midrule
\pathref{pyqcd/operator/\_gradient\_flow\_gluon\_ope.py} &
\codeurl{gradient_flow_gluon_ope} &
thin link $\to$ Wilson flow $\to$ schema-v2 OPE；保存/加载和严格验证 \\
\midrule
\pathref{pyqcd/operator/\_gradient\_flow\_gluon\_ope.py} &
\codeurl{flowed_gluon_ope} &
对已流化 gauge 计算 Clover、对偶、$\pm z$ Wilson 线和六组件 \\
\midrule
\pathref{pyqcd/analysis/\_gradient\_flow\_ratio.py} &
\codeurl{calculate_gradient_flow_ratios} &
全 channel/orientation/component 的 C3/C2 与 delete-one jackknife \\
\midrule
\pathref{pyqcd/analysis/\_gradient\_flow\_ratio.py} &
\codeurl{physical_helicity_ratio} &
保留复数 ratio，按需取 helicity orientation/component 的末端相位 \\
\midrule
\pathref{pyqcd/renorm/\_gradient\_flow\_quasi.py} &
\codeurl{match_gradient_flow_ope_to_quasi} &
参考一圈 $\GF\to\MSbar$ 准算符端点因子与组合重构 \\
\midrule
\pathref{pyqcd/pipeline/\_gradient\_flow\_gluon.py} &
\codeurl{run_gradient_flow_gluon_ope} &
ILDG \code{.lime}/\code{.lime.contents} 读取和单组态 artifact 驱动 \\
\midrule
\pathref{examples/pyqcd/gradient\_flow\_gluon\_ope.py} &
命令行入口 &
将实际 reader、流、OPE 和配对 NPY/JSON 输出串联 \\
\bottomrule
\end{tabularx}
\end{table}

\subsection{端到端数据流}

\begin{center}
\begin{tikzpicture}[
  node distance=0.33cm and 0.16cm,
  box/.style={
    draw=pyblue,fill=pyblue!7,rounded corners=2pt,
    align=center,minimum height=0.82cm,text width=2.65cm,
    font=\small
  },
  arrow/.style={-{Latex[length=2.2mm]},draw=pyblue,semithick}
]
\node[box] (u) {ILDG thin link\\$U_\mu(x)\in\SU(3)$};
\node[box,right=of u] (flow) {4D Wilson flow\\RK3, $\tau,\epsilon$};
\node[box,right=of flow] (ope) {Clover + dual\\$\pm z$ Wilson line};
\node[box,right=of ope] (ratio) {C3/C2\\pol35/nopol + JK};
\node[box,right=of ratio] (quasi) {$\GF\to\MSbar$\\准算符因子};
\draw[arrow] (u) -- (flow);
\draw[arrow] (flow) -- (ope);
\draw[arrow] (ope) -- (ratio);
\draw[arrow] (ratio) -- (quasi);
\end{tikzpicture}
\end{center}

上述最后一步只是 reference 一圈端点转换。若目标是物理 gluon TMD-PDF，
还需要在同一方案中处理矩阵元重整化、gluon--quark mixing、有限坐标重构、
软因子/快度减法、LaMET 或 pseudo-PDF matching，以及 $a\to0$ 等极限。

\section{对照测试与数值证据}

\subsection{契约测试}

新增
\pathref{pyqcd/testing/\_gradient\_flow\_gluon\_contract.py} 覆盖九个独立
测试方法：

\begin{enumerate}
  \item schema、轴标签、metadata 和零流时间；
  \item traceless projection 与公共 Clover observable；
  \item 完整空间和在随机局域规范变换下的不变性；
  \item 六组件、方向和 helicity selector 恒等式；
  \item $\epsilon=0.01$ 默认步长以及细步长收敛；
  \item 一圈 GF-to-quasi 分量因子和 combined 重构；
  \item full complex pol35 covariance 与 nopol 分母；
  \item ratio 轴、mask、末端 helicity 相位；
  \item NPY/JSON 原子 round-trip。
\end{enumerate}

\begin{table}[htbp]
\centering
\caption{本轮可复现的数值验证结果}
\begin{tabularx}{\textwidth}{@{}p{0.31\textwidth}p{0.20\textwidth}X@{}}
\toprule
验证项 & 结果 & 解释 \\
\midrule
独立 OPE/ratio 契约 &
$9/9$ &
随机 $\SU(3)$ 小格点和合成 C2/OPE 全部通过 \\
\midrule
ratio 中心 $C_3$ &
$\max|{\rm diff}|=0$ &
PyQCD 与参考 \code{calc\_ratio\_gradient\_flow.py} 逐项一致 \\
\midrule
ratio 中心 $C_2$ &
$\max|{\rm diff}|=0$ &
共同 nopol 分母一致 \\
\midrule
central ratio &
$\max|{\rm diff}|=0$ &
完整复数除法和轴布局一致 \\
\midrule
jackknife 均值/偏差修正 &
$\max|{\rm diff}|=0$ &
删除每个组态后同时重算真空扣除和分母 \\
\midrule
jackknife 实部/虚部误差 &
$\max|{\rm diff}|=0$ &
误差按复数的实、虚分量分别统计 \\
\midrule
有效 insertion mask &
完全一致 &
有效区间 $0\leq\tau_{\rm ins}\leq t_{\rm sep}$，padding 为 NaN \\
\midrule
新 OPE 与 legacy 单组件 &
约 $10^{-18}$ &
同一流化 gauge 上的单组件数值重合；不是生产数据逐元素证据 \\
\midrule
随机局域规范变换全 OPE &
约 $2.1\times10^{-16}$ &
完整空间和的双局域 trace 规范不变 \\
\midrule
$\epsilon=0.01$ 对 $\epsilon=0.005$ &
相对差约 $3.1\times10^{-6}$ &
小格点有限步长误差量级符合三阶流积分预期 \\
\bottomrule
\end{tabularx}
\end{table}

\subsection{仓库级回归}

执行命令：
\begin{lstlisting}[language=bash]
python -m pyqcd.testing._gradient_flow_gluon_contract
python examples/pyqcd/conftest.py
python -m compileall -q pyqcd examples/pyqcd
git diff --check
\end{lstlisting}

结果为：
\begin{center}
\fbox{\parbox{0.84\textwidth}{
\centering
\texttt{9 个新契约方法通过}\\
\texttt{106 passed, 0 skipped, 0 failed}\\
\texttt{compileall 通过；git diff --check 通过}
}}
\end{center}

全量回归覆盖既有 Wilson flow SU(3)/耗散、TMD 算符、HYP/Stout、torch/cupy
适配、匹配核、$Z_R$、混合方案、外推、蒸馏管线、MPI 和持久化契约。因而新增
模块没有改写旧路径默认行为。

\subsection{测试证据的正确解读}

上表的 ratio 零差异来自\textbf{同一合成/小数组输入}经两份 estimator 的
直接对照，不是从当前工作树中恢复了参考的 $231$ 组态产物。OPE 的规范变换
和 legacy 单组件测试证明几何与局域代数正确，但也不等价于真实 ensemble
统计、流时间窗口、动量依赖和拟合平台已经验证。

\section{物理正确性、局限与未闭合链}

\subsection{已能作出的物理判断}

\begin{itemize}
  \item Wilson flow 是四维、规范协变的耗散演化；正流时间降低 Wilson
        plaquette 作用量密度并保持链接接近 $\SU(3)$。
  \item Clover 场强、对偶场强和直线 Wilson 线属于同一 flowed gauge，
        双局域 trace 的规范变换在端点相消。
  \item $O_U$ 的 $M_{titi}-M_{ijij}$ 和 helicity 原始组件的方向 odd
        投影可由代码中的线性恒等式直接验收。
  \item disconnected ratio 的 vacuum subtraction 和 pol35/nopol 分工
        与参考 v2 一致，且完整复数统计没有被提前截断。
\end{itemize}

\subsection{不能从本轮结果推出的结论}

\begin{itemize}
  \item 不能把有限流 bare OPE 或 bare ratio 直接称为物理 gluon PDF；
  \item 不能用单格距、少量 $\tau$ 点宣布小流时间窗口已经受控；
  \item 不能由随机小格点的数值稳定性推出 $N_{\rm conf}=231$ 的统计误差、
        helicity 信号或 momentum dependence；
  \item 不能把已有 PyQCD 的 TMD staple、$Z_R$、hybrid、CS、SFTX 和 NLO
        入口自动视为已经消费了 schema-v2 ratio；
  \item 不能声称当前报告完成了 gluon--quark mixing、软函数/rapidity
        subtraction、LaMET matching 或连续极限。
\end{itemize}

\subsection{下一阶段的最小闭环}

要把算法级交付推进为真实物理生产，至少需要：

\begin{enumerate}
  \item 按参考共同 manifest 准备完整 $L32\times96$ thin-link 组态、C2
        \code{nopol/pol35} 和流时间列表；
  \item 对每个组态生成配对 NPY/JSON OPE artifact，并保存输入哈希、
        flow 参数、projection、源版本和失败日志；
  \item 在同一配置索引下运行 ratio，验证有效/无效 mask、共享分母和
        jackknife；
  \item 对 $\tau$、$z/a$、$P_z$、动量和 helicity 组合进行窗口敏感性，
        不能用较大流时间的较小误差替代小流控制；
  \item 明确选择 GF-to-quasi 一圈转换、$Z_R$/混合或其它同一方案的
        renormalization，并将其与 ratio metadata 绑定；
  \item 补齐软因子、快度演化、NLO kernel、有限坐标 Fourier 重构，以及
        $a,P_z,m_\pi,L$ 的联合外推；
  \item 只有在上述对象和质量门均通过后，才报告带系统误差的物理
        $xg(x,b_\perp;\mu,\zeta)$。
\end{enumerate}

\section{最终结论}

本轮完成了参考目录中最关键的 Gradient Flow/GluonPDF 算法层移植与修正：

\begin{equation}
\boxed{\begin{aligned}
\text{thin link}
&\longrightarrow\text{4D Wilson flow RK3}
 \longrightarrow\text{Clover/dual + straight OPE}\\
&\longrightarrow\text{complex disconnected ratio}
\end{aligned}}
\end{equation}

同时，参考一圈 $\GF\to\MSbar$ 准算符端点转换已作为独立 renorm 模块加入，
避免与既有 SFTX/TMD API 混淆。新增代码通过 $9$ 个专门契约方法和仓库
$106$ 项全量回归；核心 estimator 与参考脚本逐项零差异。

最重要的剩余边界是数据而非口头结论：当前工作树没有参考的完整
$L32\times96$、$231$ 组态 OPE/2pt 生产产品，所以交付等级是
\textbf{算法级完整、真实生产级待补数据}。在补齐共同 manifest、真实产物、
重整化/混合/软因子/matching 和连续极限之前，最终对象应严格称为
\textbf{有限流 bare disconnected gluon observable}，而不是已经完成的
物理 gluon PDF。

\appendix
\section{最终文件与验证命令}

\subsection{新增或修改文件}

\begin{itemize}
  \item \pathref{pyqcd/operator/\_gradient\_flow\_gluon\_ope.py}；
  \item \pathref{pyqcd/analysis/\_gradient\_flow\_ratio.py}；
  \item \pathref{pyqcd/renorm/\_gradient\_flow\_quasi.py}；
  \item \pathref{pyqcd/pipeline/\_gradient\_flow\_gluon.py}；
  \item \pathref{examples/pyqcd/gradient\_flow\_gluon\_ope.py}；
  \item \pathref{pyqcd/testing/\_gradient\_flow\_gluon\_contract.py}；
  \item 相应的 \code{\_\_init\_\_.py} 导出和
        \pathref{examples/pyqcd/conftest.py} 测试调度。
\end{itemize}

\subsection{推荐调用}

\begin{lstlisting}[language=bash]
# 单组态 thin-link ILDG -> schema-v2 OPE
python examples/pyqcd/gradient_flow_gluon_ope.py \
  --input /path/to/config.lime \
  --conf-id 4150 \
  --output /path/to/output/conf4150_tau0p500_eps0.010_zdir2 \
  --nx 32 --ny 32 --nz 32 --nt 96 \
  --tau 0.5 --epsilon 0.01 --z-count 25 --z-dir 2

# 新链契约与仓库回归
python -m pyqcd.testing._gradient_flow_gluon_contract
python examples/pyqcd/conftest.py
\end{lstlisting}

\subsection{参考与工作对象的关键路径}

\begin{longtable}{@{}p{0.42\textwidth}p{0.42\textwidth}@{}}
\caption{可追溯路径索引}\\
\toprule
参考对象 & PyQCD 对应对象 \\
\midrule
\endfirsthead
\toprule
参考对象 & PyQCD 对应对象 \\
\midrule
\endhead
\pathref{refer/donghx/Gradient\_Flow\_GluonPDF/gradient-flow-gluon-ratio.md} &
\pathref{pyqcd/operator/\_gradient\_flow\_gluon\_ope.py} \\
\midrule
\pathref{refer/donghx/Gradient\_Flow\_GluonPDF/calc\_ratio\_gradient\_flow.py} &
\pathref{pyqcd/analysis/\_gradient\_flow\_ratio.py} \\
\midrule
\pathref{refer/donghx/Gradient\_Flow\_GluonPDF/gluon\_flow\_to\_quasi.py} &
\pathref{pyqcd/renorm/\_gradient\_flow\_quasi.py} \\
\midrule
\pathref{refer/donghx/Gradient\_Flow\_GluonPDF/validate\_output\_v5.py} &
\codeurl{validate_gradient_flow_gluon_ope} \\
\midrule
\pathref{refer/donghx/Gradient\_Flow\_GluonPDF/validate\_ratio.py} &
\codeurl{validate_gradient_flow_ratio_results} \\
\bottomrule
\end{longtable}

\end{document}
